\begin{lemma}\label{lem:octic-multiple-cusp-diagonal}
Let $X/F$ be a smooth projective curve with a rational cusp $c_0$,
and let $J=\operatorname{Jac}(X)$ have rank zero over $F$.
Let $C\subseteq J(F)$ be a subgroup of order $h$ generated by
$F$-rational cuspidal divisors. Suppose that $\#J(F)$ divides
$rh$. At primes $v_i$ of good reduction, suppose that the
reduction image $C_i$ has order $h$ and is identified by reducing
the same cuspidal divisors. If $D$ is an $F$-rational effective
divisor of degree $d$, then
\[
 \bigl(r[\overline D_i-d\overline c_{0,i}]\bigr)_i
 \in\operatorname{im}\bigl(C\longrightarrow\prod_i C_i\bigr).
\]
\end{lemma}
\begin{proof}
The finite group $J(F)/C$ has order dividing $r$, so
$rJ(F)\subseteq C$. Apply this to $[D-dc_0]$ and reduce.
The order hypothesis makes each reduction on $C$ injective;
relations among the common cusp generators therefore identify
the images without any choice of invariant-factor bases.
\end{proof}

\begin{remark}\label{rem:octic-diagonal-affine-action}
Suppose a group $A$ of $F$-automorphisms preserves $C$.
On classes normalized at $c_0$, its action on the prediction
$a=r[D-dc_0]$ is
\[
 a\longmapsto g(a)+rd[g(c_0)-c_0].
\]
If all divisor sets at the preceding primes were enumerated
completely and are $A$-stable, their surviving prediction set
is stable under this affine action. One may then reduce the
first support place of each new local divisor by $A$, while
retaining every possible other support and cusp complement.
Repetitions must be included. This argument requires neither
$J(F)=C$ nor linear stability of the prediction set without
the displayed base-cusp translation.
\end{remark}
